% chktex-file 44 32 36 38 18 11 8 26

\pagestyle{DS_FP}
%\NomPrenomNote{}

\bigskip


\section*{Contexte général}

Une entreprise spécialisée dans les systèmes embarqués développe un nouveau module de transmission de données sans fil. Les équipes travaillent sur trois problématiques distinctes : la modélisation du signal de charge d’un condensateur, la gestion d’une mémoire tampon (buffer), et le contrôle qualité des modules produits.

Un ingénieur modélise la tension aux bornes d’un condensateur lors d’une phase de charge particulière par une fonction de la forme :

\[f(t)=(at+b)\cdot \e^{-kt}\qquad \text{où $a$, $b$ et $k$ sont des constantes réelles.}\]

\section*{Modélisation}

\begin{itemize}
    \item $t$ est le temps en millisecondes (ms) depuis le début de la charge du condensateur.
    \item $f(t)$ est la tension en volts (V) aux bornes du condensateur à l’instant $t$.
    \item $k$ est un coefficient de décroissance qui modélise la résistance du circuit, influençant la rapidité avec laquelle la tension atteint son maximum et commence à décroître.
\end{itemize}

\section*{PARTIE A -- Etude d'une situation particulière}

Uniquement dans cette question, on pose $a=0$, $b=5$ et $k=1$.

Dans cette question uniquement, on a donc:
\[f(t)=5\cdot \e^{-t}.\]

\begin{enumerate}[series=monDS, resume]
    \item Calculer $f(0)$, que représente-t-il physiquement dans le contexte de la charge du condensateur ?
    \begin{blocvisiblex}[height=2]
    {\quad}
    \end{blocvisiblex}
    \medskip
    \item Déterminer le comportement de $f(t)$ lorsque $t$ tend vers $+\infty$ et interpréter ce résultat physiquement.
    \begin{blocvisiblex}[height=2]
    {\quad}
    \end{blocvisiblex}
    \medskip
    \item Déterminer expérimentalement au bout de combien de temps la tension sera inférieure à $0,5$ volts. (On arrondira le résultat à la milliseconde près).
    \begin{blocvisiblex}[height=1]
    {\quad}
    \end{blocvisiblex}
\end{enumerate}

\bigskip

\section*{PARTIE B -- Cas général}

Dans cette partie, on considère que $a$, $b$ et $k$ sont des constantes réelles strictement positives, mais on ne connaît pas leurs valeurs.

Des études expérimentales ont permis de déterminer que:
\begin{itemize}
    \item la tension initiale $f(0)$ est de $5$ volts ;
    \item la tension atteint un maximum au bout de $1.5$ ms.
    \item le coefficient de décroissance $k$ est égal à $0.5$.
\end{itemize}

La fonction $f$ est donc définie pour tout $t\geq 0$ par $f(t)=(at+b)\cdot \e^{-0.5t}$.

\pagebreak

\begin{enumerate}[series=monDS, resume]
    \item A l'aide des deux conditions expérimentales et à l'aide de Géogébra, déterminer les valeurs de $a$ et $b$.
    \begin{blocvisiblex}[height=1]
    {\quad}
    \end{blocvisiblex}
    \appelprof{}
\end{enumerate}

\medskip

Dans la suite de l'exercice on considèrera la fonction $f$ définie sur $[0,+\infty[$ par $f(t)=5\cdot (2t+1)\cdot \e^{-0.5t}$.

\begin{enumerate}[series=monDS, resume]
    \item Déterminer par le calcul l'expression de la fonction dérivée $f'$ et montrer que $f'(t)=5\cdot (1-t)\cdot \e^{-0.5t}$. 
    \begin{blocvisiblex}[height=3]
    {\quad}
    \end{blocvisiblex}
    \medskip
    \item Calculer la valeur exacte puis la valeur approchée de $f(3)$ à $10^{-2}$ près.
    \begin{blocvisiblex}[height=2]
    {\quad}
    \end{blocvisiblex}
    \medskip
    \item Déterminer le tableau de variations de $f$ sur $[0,+\infty[$.
    \begin{blocvisiblex}[height=5]
    {\quad}
    \end{blocvisiblex}
\end{enumerate}

\medskip

Dans les premières millisecondes suivant la mise sous tension, un ingénieur peut utiliser la tangente à l'origine comme modèle linéaire approché de la tension. 

\begin{enumerate}[series=monDS, resume]
    \item Déterminer l'équation de la tangente à la courbe représentative de $f$ au point d'abscisse $0$.
    \begin{blocvisiblex}[height=2]
    {\quad}
    \end{blocvisiblex}
    \medskip
    \item Calculer l'écart entre la valeur de $f(3)$ et la valeur de la tangente à l'origine en $t=3$.
    \begin{blocvisiblex}[height=2]
    {\quad}
    \end{blocvisiblex}
    \medskip
\end{enumerate}

\medskip

L'énergie dissipée entre l'instant 0 et l'instant $T$ est égal à $\displaystyle \int_{0}^{T} f(t) \dt$.

\begin{enumerate}[series=monDS, resume]
    \item Déterminer une primitive $F$ de $f$ sur $[0,+\infty[$.
    \begin{blocvisiblex}[height=1]
    {\quad}
    \end{blocvisiblex}
    \medskip
    \item Calculer l'énergie dissipée pendant les 5 premières millisecondes. Donner une valeur exacte puis une valeur approchée à $10^{-2}$ près. (/\getpts{Q12})
    \begin{blocvisiblex}[height=2]
    {\quad}
    \end{blocvisiblex}
    \medskip
    \item Déterminer alors l'énergie moyenne dissipée par milliseconde pendant les 5 premières millisecondes.\begin{blocvisiblex}[height=2]
    {\quad}
    \end{blocvisiblex}
    \medskip
\end{enumerate}
    

\bigskip

\section*{PARTIE C -- Controle de qualité}

L’usine produit des modules de transmission. Le service qualité distingue deux lignes de production :

\begin{itemize}
    \item la ligne A, qui produit $60\%$ des modules et génère $3\%$ de défectueux ;
    \item la ligne B, qui produit $40\%$ des modules et génère $7\%$ de défectueux.
\end{itemize}

Un module est prélevé au hasard en sortie de chaîne.

\begin{enumerate}[series=monDS, resume]
    \item Compléter l'arbre de probabilité ci-dessous.

    \item Calculer la probabilité que le module prélevé soit défectueux.
    \item Sachant que le module prélevé est défectueux, calculer la probabilité qu'il provienne de la ligne A.
\end{enumerate}

Un technicien prélève un échantillon de 20 modules indépendamment les uns des autres. On admet que la probabilité qu’un module soit défectueux est $p=0,046$

Soit $X$ la variable aléatoire qui compte le nombre de modules défectueux dans l’échantillon.

\begin{enumerate}[series=monDS, resume]
    \item Déterminer la loi de probabilité de $X$ et en donner les paramètres.
    \item Déterminer la probabilité qu'exactement 2 modules soient défectueux dans l'échantillon.
    \item Déterminer la probabilité qu'au moins 2 modules soient défectueux dans l'échantillon.
\end{enumerate}


Le technicien souhaite que la probabilité de détecter au moins un module défectueux sur $n$ prélèvements soit supérieure à $0,99$.

\begin{enumerate}[series=monDS, resume]
    \item Montrer que cette condition s'écrit ${(1-p)}^n < 0.01$.
    \item A l'aide de Géogébra, déterminer $n$.
    \appelprof{}
\end{enumerate}

\section*{PARTIE D -- Gestion d'une mémoire tampon}

Un buffer (mémoire tampon) contient initialment 200 octets de données. A chaque cycle d'horloge, le système traite et supprime $10\%$ du contenu actuel du buffer.

On note $u_n$ la quantité de données (en octets) présente dans le buffer après $n$ cycles d'horloge.

\begin{enumerate}[series=monDS, resume]
    \item Justifier que la suite $(u_n)$ vérifie la relation de récurrence $u_{n+1}=0.9\cdot u_n$.
    \item Déterminer la nature de la suite $(u_n)$. 
    \item Exprimer $u_n$ en fonction de $n$. 
    \item Déterminer la limite de la suite $(u_n)$ lorsque $n$ tend vers $+\infty$. Interpréter ce résultat dans le contexte de la gestion du buffer.
    \item Compléter l'algorithme ci-dessous qui calcule le nombre de cycles d'horloges nécessaires pour que la quantité de données dans le buffer soit inférieure à 1 octet.

\begin{verbatim}
n = 0
u = 200
while .........:
    u = .........
    n = .........
print(n)
\end{verbatim}
    \item Par la méthode de votre choix, déterminer le nombre de cycles d'horloge nécessaires pour que la quantité de données dans le buffer soit inférieure à 1 octet.
\end{enumerate}

\bigskip



